\section{Conclusion}
\label{sec:conclusion}

We compared typed decision models with their own backbones and with trained classifiers on identical items and
candidate sets, under controlled prevalence shift and with 35 tests fixed before any test was computed. Vendor decision
training changed the same backbone in a task-dependent direction: with only an adapter between them, Nimble-9B was better
calibrated than Qwen3.5-9B on intent tasks and 12 Macro-F1~points worse on financial sentiment (RQ1). Under moderate
prevalence shift, Jev's calibration error and threshold risk grew significantly more than those of a fine-tuned DeBERTa
only on financial sentiment (by 1.7 and 2.5~points); this difference depended on the shift scenarios, and averaged over
many scenarios it was at most one point on every dataset (RQ2). A label-free EM correction with the model's implicit
prior lowered Jev's calibration error under shift by 2.7 and 1.8~points on the two datasets where it was highest,
consistently in sign across scenarios, but recalibration on labelled development data reduced it much further (RQ3).
The three typed decision models we tested are neither uniformly better than their backbones nor uniquely fragile under
prevalence shift; for calibration, their main weakness is the miscalibration they already have without shift. Whether
one is enough for a task is best decided by testing it against its base model and recalibrating it on labelled data
from the deployment setting.
